\chapter{INTRODUÇÃO}

\section{Motivação operacional}

Sistemas de classificação raramente tomam decisões a partir de uma única fonte de observação. Em aplicações reais, a entrada de um classificador costuma reunir medidas produzidas por sensores, resultados de exames, registros administrativos, variáveis contextuais, atributos derivados computacionalmente e avaliações realizadas por especialistas. Neste trabalho, cada uma dessas fontes observáveis é tratada como um \emph{canal de informação}. A expressão é deliberadamente ampla: um canal pode corresponder à leitura de um equipamento, a uma variável calculada a partir de várias leituras, a um resultado laboratorial, a um item de prontuário, a uma inspeção técnica, a uma auditoria humana ou a qualquer outra informação disponibilizada ao classificador.

A escolha do termo \emph{canal de informação} é intencional. O objeto de estudo do \CoInfoSim{} não são apenas os dados brutos que o classificador recebe, mas o conteúdo informacional que cada canal fornece para a tarefa de classificação. Em aprendizado de máquina, uma variável é útil na medida em que carrega informação sobre a variável-alvo e reduz a incerteza do classificador; um conjunto de canais, por sua vez, pode fornecer informação complementar ou redundante, influenciando o desempenho de forma cooperativa. Embora, na implementação, cada canal corresponda a um atributo ou a um subconjunto de atributos da tabela de dados, o que o projeto investiga é a relação informacional entre os canais e a variável-alvo, e também entre os próprios canais. Por isso, \emph{canal de informação} representa melhor a motivação científica do trabalho do que \emph{canal de dados}: o interesse está na informação preditiva disponibilizada ao classificador, e não apenas no dado físico que chega até ele.

A utilização de múltiplos canais pode ampliar a informação disponível para a decisão, mas também introduz custos. Esses custos não se resumem ao preço de aquisição de um sensor. Dependendo da aplicação, podem envolver implantação, calibração, manutenção, consumo de energia, transmissão, armazenamento, processamento, latência, risco, invasividade, disponibilidade de profissionais e tempo necessário para produzir a informação. A obtenção do próprio rótulo de referência também pode ser onerosa: um alvo pode depender de um exame laboratorial, de um sensor certificado mais caro, de acompanhamento longitudinal, de auditoria ou de avaliação especializada. A literatura de aprendizagem sensível a custos reconhece que aplicações reais envolvem diversas categorias de custo além do erro de classificação \cite{turney2002}.

A estrutura de custos é específica de cada domínio. Um mesmo procedimento pode produzir vários canais simultaneamente, de modo que retirar uma das variáveis resultantes não necessariamente reduz o custo da operação. Em sentido oposto, um sensor já instalado pode possuir custo fixo elevado e custo marginal relativamente baixo por nova amostra, enquanto uma análise humana pode apresentar custo predominantemente variável. Também é possível que a coleta de um canal dependa da frequência de amostragem, do tempo de funcionamento ou da infraestrutura compartilhada com outros canais. Consequentemente, este trabalho não pressupõe uma função universal de custo. A motivação operacional é anterior à modelagem econômica: antes de atribuir valores monetários, é necessário compreender quais canais são redundantes, quais são complementares e em que condições uma combinação passa a produzir vantagem.

Essa motivação acompanha a linha de pesquisa desde seus primeiros simuladores. A pergunta operacional subjacente pode ser formulada da seguinte maneira: todos os canais mantidos por uma operação contribuem de forma relevante para a classificação, ou existem configurações menores que preservam o resultado com menor complexidade? A resposta não pode ser obtida apenas pela avaliação isolada de cada variável. Como discutem \textcite{guyon2003}, a seleção de variáveis pode servir não somente para melhorar a predição, mas também para reduzir requisitos de medição e processamento e tornar o modelo mais compreensível. Entretanto, uma variável aparentemente fraca pode se tornar valiosa quando combinada com outra, enquanto duas variáveis individualmente fortes podem carregar informação redundante.

O planejamento de uma arquitetura informacional exige, portanto, uma análise conjunta. Não basta perguntar qual é o melhor canal isolado. É necessário comparar subconjuntos, observar como eles se ordenam e identificar quando subconjuntos maiores passam a justificar sua complexidade. Essa transição pode depender da quantidade de exemplos rotulados disponíveis: uma combinação de canais pode conter mais informação potencial, mas exigir mais dados para que o classificador consiga estimar sua estrutura de maneira estável.

\section{Problema científico}

Considere uma tarefa supervisionada com $d$ canais candidatos. Para cada subconjunto de canais, é possível treinar um classificador e avaliar sua perda. Em princípio, o subconjunto com mais canais possui acesso a toda a informação disponível aos subconjuntos menores. Na prática, porém, a inclusão de novas dimensões aumenta o número de relações que precisam ser estimadas e pode elevar a variância do processo de aprendizagem. Em regimes de pequenas amostras, um canal adicional pode não produzir vantagem ou até deteriorar o desempenho. O fenômeno clássico discutido por \textcite{hughes1968} e aprofundado por \textcite{raudys1991} mostra que a relação entre dimensionalidade, complexidade e tamanho amostral é central na classificação estatística.

Desse modo, a cooperação não é uma propriedade fixa de um conjunto de variáveis. Ela depende, pelo menos, de quatro elementos: o subconjunto de canais observado, o modelo de classificação, a distribuição conjunta dos dados e o tamanho amostral utilizado no treinamento. Um subconjunto pode ser superior em determinado regime e inferior em outro. Dois classificadores também podem responder de maneira distinta às mesmas relações entre canais, porque impõem hipóteses diferentes sobre fronteiras de decisão, dependências e interações.

O estudo dessas relações por meio de dados reais encontra uma limitação prática. Para observar como a cooperação evolui em uma ampla faixa de tamanhos amostrais, é preciso repetir o treinamento muitas vezes, controlar a composição das amostras e, idealmente, investigar regimes com mais dados do que os atualmente disponíveis. Modelos sintéticos oferecem um ambiente controlado para essa análise, mas sua utilidade não deve ser presumida apenas porque produzem pontos visualmente semelhantes aos dados reais ou porque sustentam uma perda próxima em uma configuração particular.

O problema científico deste relatório consiste em comparar a \emph{estrutura de cooperação} que emerge quando classificadores são treinados com dados reais com aquela que emerge quando os mesmos classificadores são treinados com amostras sintéticas produzidas por modelos parametrizados a partir dos dados reais. Essa estrutura inclui a ordenação dos subconjuntos, as relações de vitória entre pares e os limiares amostrais nos quais uma combinação passa a superar outra. Portanto, o objeto principal não é a fidelidade visual das amostras sintéticas nem a reprodução exata da perda absoluta. O interesse recai sobre a preservação das relações que orientam a comparação entre configurações de canais.

Essa distinção é especialmente importante porque fidelidade geométrica, fidelidade preditiva e fidelidade estrutural não são equivalentes. Um modelo gerador pode produzir distribuições visualmente semelhantes às dos dados reais e, ainda assim, alterar o ranking dos subconjuntos. Em sentido inverso, uma representação sintética aparentemente grosseira pode preservar relações decisórias relevantes. Também é possível que um classificador com desempenho preditivo absoluto baixo responda de maneira semelhante às vantagens relativas entre combinações de canais quando treinado com dados reais e quando treinado com dados sintéticos, desde que ambos os treinamentos sejam avaliados sobre o mesmo conjunto de teste real. Os experimentos discutidos neste relatório foram desenhados para tornar essas diferenças observáveis.

\section{Evolução do \SLACGS{} ao \CoInfoSim}

O \CoInfoSim{} é uma evolução de uma linha anterior de simuladores voltados ao estudo da cooperação em classificação. O ponto de partida foi o \SLACGS{}, que utilizava distribuições gaussianas especificadas de forma controlada e simulação de Monte Carlo para acompanhar curvas de perda em diferentes tamanhos amostrais. O sistema anterior consolidou mecanismos que permanecem importantes: geração reprodutível, crescimento incremental das amostras, repetição adaptativa e relatórios que agregam resultados de múltiplas simulações \cite{azevedo2026coinfosim}.

No \SLACGS{}, a comparação era organizada principalmente por crescimento incremental de dimensionalidade. Partia-se de uma dimensão, acrescentava-se uma segunda, depois uma terceira, formando uma cadeia prefixada de modelos. Esse desenho permitia estudar quando a inclusão de uma nova dimensão produzia vantagem, mas não avaliava todas as combinações possíveis. Se a ordem dos canais fosse $X_1,X_2,X_3$, por exemplo, a sequência podia comparar $\{X_1\}$, $\{X_1,X_2\}$ e $\{X_1,X_2,X_3\}$, sem necessariamente investigar $\{X_2\}$, $\{X_3\}$ ou $\{X_1,X_3\}$ como objetos independentes.

A primeira reformulação do \CoInfoSim{} consiste em substituir essa cadeia incremental pela enumeração dos subconjuntos não vazios. O objeto deixa de ser apenas a dimensionalidade e passa a ser a composição informacional. Essa mudança permite distinguir, por exemplo, uma dupla formada por canais fortes mas redundantes de outra dupla composta por canais individualmente modestos e complementares. Também permite acompanhar a mudança de ranking entre subconjuntos de mesma cardinalidade, algo que a comparação puramente dimensional não revela.

A segunda reformulação consiste em ancorar os experimentos em dados reais. O trabalho anterior foi útil para isolar mecanismos em cenários gaussianos controlados, mas sua validade externa era limitada. Durante a avaliação daquela linha, foi apontada a ausência de exemplos empíricos capazes de mostrar se a estrutura observada em modelos idealizados também surgia em problemas reais. A crítica não invalida o valor dos cenários controlados, mas delimita seu alcance: um resultado demonstrado sob parâmetros escolhidos manualmente não basta para sustentar conclusões sobre aplicações concretas.

A disciplina de Ciência de Dados forneceu o contexto para essa transição. A incorporação de datasets públicos exigiu decisões que não apareciam no simulador idealizado: interpretação do domínio, definição de alvos, prevenção de vazamento de informação, escolha de particionamento, tratamento de valores ausentes, padronização, balanceamento e análise exploratória da geometria dos dados. Alguns datasets inicialmente considerados foram adiados precisamente porque sua estrutura temporal ou hierárquica não era compatível com um tratamento tabular ingênuo. Assim, o desenvolvimento do projeto passou a incluir não apenas a execução de classificadores, mas a construção de protocolos científicos específicos para cada conjunto de dados.

O protocolo resultante compara três condições experimentais de treinamento, adiante denominadas \emph{braços}. No primeiro braço, amostras balanceadas são retiradas do conjunto de treinamento real e avaliadas em um conjunto de teste composto por dados reais. Nos demais, modelos probabilísticos parametrizados somente a partir do conjunto de treinamento real produzem as amostras sintéticas usadas para treinar os classificadores; a avaliação também ocorre no mesmo conjunto de teste real fixo. Os geradores atuais são uma Gaussiana única por classe e uma mistura gaussiana condicional à classe. Esse desenho se relaciona à lógica \emph{Train on Synthetic, Test on Real}, proposta na literatura de avaliação de dados sintéticos \cite{esteban2017}, mas desloca a ênfase da utilidade preditiva global para a preservação da organização cooperativa que emerge do treinamento dos classificadores para os diferentes subconjuntos.

A ampliação de escopo também motivou a mudança terminológica de sensores para canais de informação. Sensores permanecem uma aplicação importante, especialmente nos cenários de ocupação e qualidade do ar. Entretanto, o SUPPORT2 mostra que a mesma formulação pode representar medidas fisiológicas e laboratoriais. Em aplicações futuras, os canais podem incluir informações produzidas por pessoas ou por sistemas computacionais. O nome \CoInfoSim{} reflete essa generalização.

\section{Pergunta central e objetivos}

A pergunta central deste trabalho é:

\begin{researchquestionbox}[Pergunta central]
Em que medida a estrutura de cooperação entre canais de informação que emerge do treinamento de classificadores com amostras sintéticas, produzidas por modelos parametrizados a partir de dados reais, reproduz aquela que emerge do treinamento dos mesmos classificadores com amostras reais?
\end{researchquestionbox}

O objetivo geral é comparar a estrutura de cooperação que emerge quando os classificadores são treinados com amostras reais com aquela que emerge quando os mesmos classificadores são treinados com amostras sintéticas, produzidas por uma Gaussiana única ou por uma mistura de gaussianas parametrizadas a partir dos dados reais de treinamento. Em todas as condições experimentais, o desempenho é avaliado sobre o mesmo conjunto de teste composto por dados reais.

Para atingir esse objetivo, foram definidos os seguintes objetivos específicos:

\begin{enumerate}
  \item formular o problema em termos de canais de informação e de seus subconjuntos;
  \item construir protocolos reprodutíveis para datasets públicos de domínios distintos;
  \item avaliar exaustivamente os subconjuntos não vazios dos canais selecionados;
  \item acompanhar a perda e a organização dos subconjuntos ao longo de uma grade crescente de tamanhos amostrais por classe;
  \item comparar os braços Real$\rightarrow$Real, Gaussiana$\rightarrow$Real e GMM$\rightarrow$Real;
  \item avaliar a dependência dos resultados em relação ao classificador;
  \item medir separadamente a preservação do ranking, das relações de vitória e dos limiares $\Nstar$;
  \item identificar situações em que uma representação sintética simples é estruturalmente suficiente;
  \item discutir como os objetos produzidos pelo simulador podem subsidiar, em trabalhos futuros, planejamento operacional e modelagem de custos.
\end{enumerate}

\section{Contribuições do trabalho}

As contribuições deste relatório são metodológicas, experimentais e de engenharia científica. No plano conceitual, o trabalho generaliza a cooperação entre sensores para cooperação entre canais de informação e substitui a comparação prefixada de dimensões pela avaliação dos subconjuntos. No plano experimental, introduz um protocolo de três braços em que todos os classificadores são avaliados sobre um teste real fixo, permitindo comparar a estrutura produzida por treinamento real e sintético.

No plano das métricas, a análise não se limita ao melhor resultado final. O sistema produz rankings para cada tamanho amostral, compara matrizes de relações par a par por meio de \WinnersAgreement{} e avalia a similaridade progressiva das estruturas de $\Nstar$. Essas medidas são mantidas separadas porque representam componentes distintos da organização cooperativa. A ausência de um índice agregado evita ocultar divergências relevantes entre ordenação, relações de vitória e posição dos cruzamentos.

No plano empírico, são estudados três datasets públicos: Occupancy Detection, UCI Air Quality e SUPPORT2 com mortalidade em 180 dias. Os conjuntos representam ocupação de ambiente, monitoramento ambiental e prognóstico clínico, respectivamente. O SUPPORT2 também recebe um experimento adicional com Random Forest, permitindo observar a preservação estrutural em um classificador não linear e computacionalmente mais oneroso.

Finalmente, o projeto consolida uma infraestrutura de simulação de Monte Carlo com amostras incrementalmente aninhadas, critério adaptativo de parada, paralelização, persistência de metadados e geração automatizada de relatórios. Essa infraestrutura não é apenas um produto de software: ela é parte do controle experimental necessário para que as comparações sejam reproduzíveis e auditáveis.

\section{Organização do relatório}

O Capítulo 2 apresenta os conceitos de canais, subconjuntos, cooperação, pequenas amostras, dados sintéticos, classificadores e custos. O Capítulo 3 formaliza as perguntas de pesquisa e as hipóteses. O Capítulo 4 descreve a seleção e a preparação dos datasets. O Capítulo 5 detalha o protocolo experimental e a infraestrutura de Monte Carlo. O Capítulo 6 define os objetos estruturais e as métricas de preservação. O Capítulo 7 apresenta os resultados dos três cenários. O Capítulo 8 realiza a discussão transversal. O Capítulo 9 retoma a motivação operacional e delimita a agenda futura de custos. O Capítulo 10 discute limitações e ameaças à validade, e o Capítulo 11 apresenta as conclusões.
